\scanpage{445}
\mspage{525}

\uline{The Center of} $\mathcal M_{\rho,\sigma}$. It maps under $\widetilde\Theta$ into the center of $PG$ which is trivial, hence
\[
\uline{\mathrm{Cent}}(\mathcal M_{\rho,\sigma})\subset\operatorname{Ker}\widetilde\Theta
=R_{\rho,\sigma}\simeq S'T_\rho\times SN'_\sigma\times G_m
\]
\unsure{whence} obviously. Since $R_{\rho,\sigma}$ commutes with $SG$, and $R_{\rho,\sigma}\,\deleted{\unsure{\times}}\cdot SG=\mathcal M_{\rho,\sigma}$, we conclude
\[
\tag{128}
\uline{\mathrm{Cent}}(\mathcal M_{\rho,\sigma})
=\uline{\mathrm{Cent}}(R_{\rho,\sigma})
=S'T_\rho\times\underset{?}{\uline{\mathrm{Cent}}}(SN'_\sigma)\times G_m
\]
let us therefore determine $\mathrm{Cent}(SN'_\sigma)$. Let us note that
\[
\left\{\begin{array}{l}
\mu_2\simeq N'_\sigma/T_\sigma\quad\text{operates faithfully on }T_\sigma\\
\text{and }T_\sigma=G_m.ST_\sigma
\end{array}\right.
\qquad
\begin{array}{|l}
\text{hence }\underset{\substack{\rotatebox{90}{$\simeq$}\\SN'_\sigma/\unsure{T_\sigma}}}{\mu_2}\text{ operates faithfully}\\
\text{on }ST_\sigma
\end{array}
\]
whence
\[
\tag{129}
\uline{\mathrm{Cent}}(SN'_\sigma)\subset ST_\sigma
\quad\text{whence}\quad
\uline{\mathrm{Cent}}(SN'_\sigma)\simeq(ST_\sigma)^{\mu_2}.
\]

Now we find that
\inserted{\[
\tag{130}
T_\sigma^{\mu_2}\simeq G_m
\]}
\deleted{\unsure{the section} $-1$ of $\mu_2$ operates on $ST_\sigma$ by $x\longrightarrow x^{-1}$, which shows that
\[
\uline{\mathrm{Cent}}(SN'_\sigma)\subset {}_2(ST_\sigma)
=\left\{\unsure{c}\sigma+d\ \middle|\ \begin{array}{l}
-\nu c^2+d^2=1\\
\nu c^2+d^2+2\sigma cd=0
\end{array}\right\}
\]
We have for a section of ${}_2ST_\sigma$, \deleted{\unsure{g|that}}
\[
\begin{array}{ll}
g=c\sigma+d & \det g=-\nu c^2+d^2=1\\
&g^2=\nu c^2+d^2+2cd\sigma=1
\end{array}
\]
\[
\text{i.e.}\quad
\left\{\begin{array}{ll}
-\nu c^2+d^2=1 & \nu c^2+d^2=1\\
2cd=0 &
\end{array}\right.
\]
\deleted{whence} i.e.
\[
\left\{\begin{array}{ll}
2\nu c^2=0 & \text{hence }2c^2=0\quad(\text{since }\nu\text{ invertible}).\\
2cd=0 &\\
-\nu c^2+d^2=1 &
\end{array}\right.
\]
if $2$ invertible this amounts to
\[
\left\{\begin{array}{lll}
&c^2=0&\\
\deleted{\unsure{c^2=0}}&d^2=1&(\text{hence }d\text{ invertible})\\
\deleted{\unsure{d^2=1}}&\unsure{c}d=0&\text{hence }c=0
\end{array}\right.
\]
\unsure{hence}
\[
\boxed{c=0,\ d^2=1}\quad\text{i.e. }g\in\mu_2
\]
\unsure{for|hence if} $2$ invertible.}
\inserted{whence}
\[
\tag{\unsure{131|129}}
\uline{\mathrm{Cent}}(SN'_\sigma)
=\unsure{ST_\sigma}^{\mu_2}
=ST_\sigma\cap G_m
\inserted{=\mu_2}\quad
\deleted{\text{if }2\text{ invertible on }S.}
\]

\deleted{In the general case we will only have
\[
\tag{130}
\mu_2\subset\uline{\mathrm{Cent}}(SN'_\sigma)
\]
and this inclusion induces an isomorphism on the groups of sections in the case where $2$ is a \unsure{regular} section of}

\mspage{444}

\begin{center}
{\small
444\\
\texttt{https://grothendieck.umontpellier.fr}\\
Université de Montpellier, UMR CNRS 5149, Case courrier 051, Place Eugène Bataillon\\
34095 Montpellier cedex 5 - France
}
\end{center}

\scanpage{446}

\deleted{$\uline{\mathcal O}_S$. Let us denote by \deleted{$\unsure{\mu'_2}$}
\[
\tag{131}
\begin{gathered}
\mu'_2\subset ST_\sigma\\[-.5em]
\rotatebox{90}{$=$}\\[-.5em]
\deleted{\unsure{\uline{\mathrm{Cent}}}}(ST_\sigma)^{\mu_2}
\end{gathered}
\]
\unsure{this|a} subgroup of}

whence \unsure{finally|finally}, \unsure{namely|i.e.}:
\[
\tag{128}
\underset{\unsure{=}}{\uline{\mathrm{Cent}}(\mathcal M_{\rho,\sigma})}
\simeq S'T_\rho\times\mu_2\times G_m
\subset R_{\rho,\sigma}\simeq S'T_\rho\times SN'_\sigma\times G_m
\]
whence
\[
\tag{129}
\uline{\mathrm{Cent}}(\mathcal M_{\rho,\sigma})\cap\mathcal M_{\rho,\sigma}^{\circ}
=ST_\rho\times\mu_2\times G_m
\subset R_{\rho,\sigma}^{\circ}\simeq ST_\rho\times ST_\sigma\times G_m
\]

Let us return to the integral sections in the absolute modular case
\[
\tag{130}
R_{\rho,\sigma}(\mathbb Z)\simeq
\underbrace{S'T_\rho(\mathbb Z)}_{\unsure{T_\rho}(\mathbb Z)\simeq\mathbb Z/6\mathbb Z}
\times
\underbrace{SN'_\sigma(\mathbb Z)}_{T_\sigma(\mathbb Z)=\mathbb Z/4\mathbb Z}
\times
\underbrace{G_m(\mathbb Z)}_{\mu_2(\mathbb Z)\simeq\pm1}
\]
\[
\tag{131}
\begin{aligned}
\uline{\mathrm{Cent}}(\mathcal M_{\rho,\sigma})\unsure{[\mathbb Z]|(\mathbb Z)}
&\simeq S'(T_\rho\unsure{[\mathbb Z]|(\mathbb Z)})\times\mu_2(\mathbb Z)\times G_m(\mathbb Z)\\
&\simeq\mathbb Z/6\mathbb Z\times\{\pm1\}\times\{\pm1\}
\end{aligned}
\]
and for the record (cf. (75), p. 514)
\[
\tag{132}
\begin{aligned}
\underbrace{\uline{\mathrm{Cent}}(\mathcal M_{\rho,\sigma}^{\circ})}_{\substack{R_{\rho,\sigma}^{\circ}\\\simeq ST_\rho\times ST_\sigma\times G_m}}(\mathbb Z)
&\simeq ST_\rho(\mathbb Z)\times ST_\sigma(\mathbb Z)\times G_m(\mathbb Z)\\
&\simeq\mathbb Z/6\mathbb Z\times\mathbb Z/4\mathbb Z\times\{\pm1\}.
\end{aligned}
\]

In summary, we have the diagram of \inserted{strict} inclusions of group schemes
\[
\tag{133}
\diagram{\begin{tikzcd}[ampersand replacement=\&,column sep=1.5em,row sep=5em]
\text{dimension }3 \&
\begin{gathered}
\uline{\mathrm{Cent}}(\mathcal M_{\rho,\sigma}^{\circ})=R_{\rho,\sigma}^{\circ}\\
\simeq ST_\rho\times ST_\sigma\times G_m
\end{gathered}
\arrow[rr,hook,"{{{\substack{\text{(quotient}\\\mu_2\times\mu_2\text{)}}}}}"']
\& \&
\begin{gathered}
R_{\rho,\sigma}\\
\simeq S'T_\rho\times SN'_\sigma\times G_m
\end{gathered}
\&\text{dimension }3\\
\text{dimension }2\&
\begin{gathered}
\uline{\mathrm{Cent}}(\mathcal M_{\rho,\sigma})\cap\mathcal M_{\rho,\sigma}^{\circ}\\
=ST_\rho\times\mu_2\times G_m
\end{gathered}
\arrow[u,hook]\arrow[rr,hook,"{{{\substack{\text{(quotient}\\\mu_2\text{)}}}}}"']
\&\&
\begin{gathered}
\uline{\mathrm{Cent}}(\mathcal M_{\rho,\sigma})\\
\simeq S'T_\rho\times\deleted{\unsure{SN'_\sigma}}\inserted{\unsure{\mu_2}}\times G_m)
\end{gathered}
\arrow[u,hook]
\&\text{dimension }2
\end{tikzcd}}
\]

\mspage{445}

\begin{center}
{\small
445\\
\texttt{https://grothendieck.umontpellier.fr}\\
Université de Montpellier, UMR CNRS 5149, Case courrier 051, Place Eugène Bataillon\\
34095 Montpellier cedex 5 - France
}
\end{center}

\scanpage{447}
\mspage{526}

valid in the general case -- and which, in the absolute modular case, induces \unsure{on} the integral points isomorphisms for the two horizontal arrows, \unsure{corresponding} resp. to the two groups
\[
\underset{\unsure{\text{and}}}{(134)}\quad
\left\{
\begin{aligned}
\uline{\mathrm{Cent}}(\mathcal{M}_{\rho,\sigma}^{\circ})(\mathbb{Z})
 &\simeq \mathbb{Z}/6\mathbb{Z}\times\mathbb{Z}/4\mathbb{Z}\times\{\pm1\}\qquad\text{and}\\
\uline{\mathrm{Cent}}(\mathcal{M}_{\rho,\sigma})(\mathbb{Z})
 &\simeq \mathbb{Z}/6\mathbb{Z}\times\{\pm1\}\times\{\pm1\}
\end{aligned}\right.
\]
which is a subgroup of it of index 2.

Finally, for the record \unsure{let us point out} once again the group
\[
(135)\qquad \deleted{\unsure{\pm1\times}}\qquad
\Sigma=\{\pm1\}\times\{\pm1\}\times\{\pm1\}
\]
contained in $\uline{\mathrm{Cent}}(\mathcal{M}_{\rho,\sigma})(\mathbb{Z})$ (of order 8, of index 3), which plays a special role in the arithmetic application, since we will have a canonical homomorphism
\[
(136)\qquad
\boxed{\diagram{\begin{tikzcd}[ampersand replacement=\&,column sep=large]
\mathcal{M}_{0,3}^{\sim}\arrow[r] \& \mathcal{M}_{\rho,\sigma}(\widehat{\mathbb{Z}})/\Sigma
\end{tikzcd}}}\,.
\]

\mspage{446}

\begin{center}
{\small
446\\
\texttt{https://grothendieck.umontpellier.fr}\\
Université de Montpellier, UMR CNRS 5149, Case courrier 051, Place Eugène Bataillon\\
34095 Montpellier cedex 5 - France
}
\end{center}

\scanpage{448}

\margin{\uline{\unsure{Where we define}}\quad\uline{$\mathcal{M}_{\rho,\sigma}$}}

\textcircled{$7^\circ$}\quad Let us make a \unsure{list|little} \unsure{recapitulation|recapitulatory} of the main group homomorphisms considered in the course of section $6^\circ$, involving the group $\mathcal{M}_{\rho,\sigma}$
\[
(134)\qquad
\diagram{\begin{tikzcd}[ampersand replacement=\&,row sep=small,column sep=large]
 \& SG\arrow[dddr,"{j_G}"] \& \& G\arrow[r]\arrow[dr,"{\delta_G=\det_G}"] \& PG\\
 \& \& \& \& G_m\\
 \& S'T_\rho\arrow[dr,"{j_\rho}"] \& \& T_\rho\arrow[r,"{\delta_\rho=\det_{T_\rho}}"] \& G_m\\
\mu_2\arrow[uuur]\arrow[ur]\arrow[dr]\arrow[dddr]
 \& \& \mathcal{M}_{\rho,\sigma}\arrow[uuur,"{\Theta=\Theta_G}"]\arrow[ur,"{\Theta_\rho}"]\arrow[dr,"{\Theta_\sigma}"]\arrow[dddr,"{\varepsilon_3}"'] \& \& \\
 \& SN'_\sigma\arrow[ur,"{j_\sigma}"'] \& \& N'_\sigma\arrow[r,"{\delta'_\sigma=\det_{T_\sigma}}"]\arrow[dr,"{\varepsilon_\sigma}"] \& G_m\\
 \& \& \& \& \mu_2\\
 \& G_m\arrow[uuur,hook,"{j_\mu}"'] \& \& \mu_2 \&
\end{tikzcd}}
\]

Let us note that this diagram can be reconstructed (with the essential properties which we are going to \unsure{pass} in \unsure{review}) from the data
\[
(135)\quad
\left\{\begin{gathered}
\diagram{\begin{tikzcd}[ampersand replacement=\&,column sep=large]
G_m\arrow[r,"{i_\rho}"] \& T_\rho\arrow[r,"{\delta_\rho}"] \& G_m \& (\delta_\rho i_\rho(\lambda)=\lambda^2)
\end{tikzcd}}\\[1em]
\diagram{\begin{tikzcd}[ampersand replacement=\&,column sep=large]
G_m\arrow[r,"{i_\sigma}"] \& N'_\sigma\arrow[r,"{\delta_\sigma}"] \& G_m \& (\delta_\sigma i_\sigma(\lambda)=\lambda^2)
\end{tikzcd}}\\[1em]
\diagram{\begin{tikzcd}[ampersand replacement=\&,column sep=large]
 \& N'_\sigma\arrow[r,"{\varepsilon_\sigma}"] \& \mu_2 \& \varepsilon_\sigma i_\sigma(\lambda)=1
\end{tikzcd}}\\[1em]
\diagram{\deleted{\begin{tikzcd}[ampersand replacement=\&,column sep=large]
1\arrow[r] \& G_m\arrow[r]\arrow[dr,phantom,"{\text{cartesian}}" description] \& G\arrow[r] \& PG\arrow[r] \& 1\\
1\arrow[r] \& \mu_2\arrow[u,hook]\arrow[r] \& SG\arrow[u,hook]\arrow[r] \& PG\arrow[u,equal]\arrow[r] \& 1
\end{tikzcd}}}
\end{gathered}\right.
\]
in addition to the \unsure{familiar} \unsure{datum}
\[
(136)\qquad
\diagram{\begin{tikzcd}[ampersand replacement=\&,column sep=large,row sep=large]
 \& 1 \& 1 \& \&\\
 \& G_m\arrow[u]\arrow[r,equal] \& G_m\arrow[u] \& \&\\
1\arrow[r] \& G_m\arrow[u,"{\lambda\mapsto\lambda^2}"]\arrow[r,hook,"{\unsure{j}}"] \& G\arrow[u,"{\delta_G}"]\arrow[r] \& PG\arrow[r] \& 1\\
1\arrow[r] \& \mu_2\arrow[u,hook]\arrow[r] \& SG\arrow[u,hook]\arrow[r] \& PG\arrow[u,equal]\arrow[r] \& 1\\
 \& 1\arrow[u] \& \unsure{0|1}\arrow[u] \& \&
\end{tikzcd}}
\]

\mspage{447}

\begin{center}
{\small
447\\
\texttt{https://grothendieck.umontpellier.fr}\\
Université de Montpellier, UMR CNRS 5149, Case courrier 051, Place Eugène Bataillon\\
34095 Montpellier cedex 5 - France
}
\end{center}

\scanpage{449}
\mspage{527}

Let us recall that we define, in terms of (135)
\[
(137)\quad
\left\{\begin{gathered}
ST_\rho=\mathrm{Ker}\,\delta_\rho,\quad S'T_\rho=\mathrm{Ker}\,\delta_\rho^2\\
T_\sigma=\mathrm{Ker}\,\varepsilon_\sigma,\quad
SN'_\sigma=\mathrm{Ker}\,\delta'_\sigma,\quad
ST_\sigma=SN'_\sigma\cap T_\sigma=\mathrm{Ker}\,\delta_\sigma\\
\hfill(\delta_\sigma=\delta'_\sigma|T_\sigma)\\[.5em]
\diagram{\begin{tikzcd}[ampersand replacement=\&,column sep=small]
\mu_2\arrow[r] \& ST_\rho\subset S'T_\rho \& \text{as composite} \& \mu_2\subset G_m\arrow[r,"{i_\rho}"] \& T_\sigma\\
\mu_2\arrow[r] \& ST_\sigma\subset SN'_\sigma \& \rule{2cm}{.3pt} \& \mu_2\subset G_m\arrow[r,"{i_\sigma}"] \& T_\sigma
\end{tikzcd}}
\end{gathered}\right.,
\]
which determines the left-hand part of diagram (134), where $\mathcal{M}_{\rho,\sigma}$ can be considered (and defined) as \inserted{\unsure{quotient of the}} \deleted{amalgamated} sum of the four factors $SG,S'T_\rho,SN'_\sigma,G_m$, \inserted{by $\mu_2$ \unsure{embedding} itself in it diagonally.}

This makes it possible in principle to reduce the description of the four fundamental morphisms from $\mathcal{M}_{\rho,\sigma}$ in (134) -- $\Theta_G,\Theta_\rho,\Theta_\sigma$ and $\varepsilon_{\unsure{3|2}}$ -- to their composites with the four homomorphisms
$j_G,j_\rho,j_\sigma,j_\mu$ -- subject to checking \inserted{that the}\deleted{\unsure{images}} of the four composites \deleted{\unsure{coincide}}\inserted{\unsure{commute}, and that their product of} \deleted{\unsure{that the|on the}}\inserted{their values on the images of an $\deleted{\unsure{\varepsilon|\equiv}}\ \varepsilon\in\Gamma\mu_2$ is trivial.} \deleted{four subgroups $\mu_2$} We will have
\[
(138)\quad
\left\{\begin{gathered}
\diagram{\begin{tikzcd}[ampersand replacement=\&,column sep=large]
\Theta_Gj_G: \& SG\arrow[r,hook] \& G,\quad\text{canonical inclusion}\quad(136)
\end{tikzcd}}\\[.5em]
\diagram{\deleted{\begin{tikzcd}[ampersand replacement=\&,column sep=small]
\unsure{\Theta_Gj_\rho}: \& \unsure{S'T_\rho}\arrow[r,hook] \& \unsure{T_\rho}\arrow[r,"{\delta_\rho}"] \& G_m\arrow[r,hook,"{\unsure{j}}"] \& G
\end{tikzcd}}}\\[.5em]
\Theta_Gj_\rho,\ \Theta_Gj_\sigma\qquad\text{trivial}\\[.5em]
\diagram{\begin{tikzcd}[ampersand replacement=\&,column sep=large]
\Theta_Gj_\mu: \& G_m\arrow[r] \& G_m\quad\text{canonical inclusion}
\end{tikzcd}}
\end{gathered}\right.
\]
(conditions of \unsure{compatibility|compatibilities} satisfied !)

\mspage{448}

\begin{center}
{\small
448\\
\texttt{https://grothendieck.umontpellier.fr}\\
Université de Montpellier, UMR CNRS 5149, Case courrier 051, Place Eugène Bataillon\\
34095 Montpellier cedex 5 - France
}
\end{center}